% Additional proved CRT variant; companion fragment to core.tex.
\section{A finite-place and bounded-size test does not ensure integrality}
The supplied real separation estimates \cite[\S1.2, equations (2)--(4)]{Input}
are necessary at an integral witness. The next theorem checks their logical
strength on explicit targets. It is not a new assertion about the Brauer group.
For the distinction between natural solutions and arithmetic approximation see
\cite[Theorems 1.1, 1.8 and 1.9]{BL}.

\begin{theorem}[Explicit rational controls for any fixed finite place set]
\label{thm:crtcontrols}
Let $\Sigma$ be any fixed finite set of rational primes. There are infinitely
many primes $p\equiv1\pmod{840}$ with a positive rational triple $(x,y,z)$ such
that the following all hold:
\begin{enumerate}
\item $4/p=1/x+1/y+1/z$, with $p/4<x<p/2$ and four separated literal roots;
\item $x,y,z\in\Z_q$ for every $q\in\Sigma\cup\{p\}$;
\item the literal prime-local completion has precisely the Type-II signatures,
Smith exponents, conductor and four-rational-point case proved above;
\item with $s_p=(p^2+3p)/4$, $B_p=p+s_p(s_p+1)$,
\[
 S<B_p,\quad\max(x,y,z)<s_p(s_p+1)/3,\quad
 \Disc(\mathcal H)\ge144/B_p^6,\quad
 2/S\le|\mathcal H'(\ell_i)|\le S^2;
\]
\item one specified prime $l\notin\Sigma\cup\{p\}$ has $v_l(z)=-1$, so this
particular triple is not integral;
\item the same $p$ has a separately displayed original integer E solution.
\end{enumerate}
\end{theorem}
\begin{proof}
Choose a prime $l\ge13$ outside $\Sigma$, and let $M=l^2\prod_{q\in\Sigma}q$.
Use the reduced progression
\[
 p\equiv1\pmod{840},\qquad p\equiv1\pmod l,\qquad p\equiv-1\pmod{11}.
                                                               \tag{IC32}
\]
Its moduli are pairwise coprime, and its residues are units. Dirichlet's theorem
\cite[Theorem 18.1]{Dirichlet} supplies infinitely many prime members. Take
$p>\max\Sigma$ and $p^2\ge128M$, and set $a=(p+3)/4$. Select $b$ by CRT:
\[
 b\equiv3\pmod p,\qquad
 b\equiv\begin{cases}1\pmod q,&q\mid a,\\0\pmod q,&q\nmid a,\end{cases}
 \quad(q\in\Sigma),
 \qquad 3b-a\equiv l\pmod{l^2}.                         \tag{IC33}
\]
There is one residue modulo $pM$. Add $pM$ to its least nonnegative representative,
so $pM\le b<2pM$. Put
\[
 c=\frac{ab}{3b-a},\qquad (x,y,z)=(a,pb,pc).             \tag{IC34}
\]
Then $a/3<c<a/2$, and direct addition proves the ES identity. For $q\in\Sigma$,
$3b-a$ is a unit: if $q\mid a$ then $q\ne3$ and $b=1$; otherwise use $b=0$.
At $l$, $a\equiv1$, so $b\equiv1/3$ and $v_l(3b-a)=1$, while $ab$ is a unit.
Thus $v_l(z)=-1$, with no cancellation.

At $p$, the normalized cluster roots are $1,3,3/11$. The character of $11$
is $-1$ by (IC32) and reciprocity. The derivative residues in (IC14) consequently
have signs $(-1,-1,1)$, because $3,5,-1$ are squares at a hard prime.
All distinctness and unit conditions used in the cluster and normalization theorems (Theorems 2.2 and 4.1) of the companion workbench proof
therefore hold literally for this rational triple, giving the four-point M case.
No global square-divisor state is inferred.

For the quantitative bounds, $p<a/3\cdot p<pc<p a/2<pb$ and $a<p$; more precisely
$pc-p>p(p-9)/12>1$ for $p\ge13$. The other two successive gaps in the ordered
list $a,p,pc,pb$ are also greater than one. Further,
\[
 S<(2M+1)p^2\le3Mp^2,
 \qquad \max(x,y,z)<2Mp^2.
\]
Since $B_p\ge p^4/16$ and $s_p(s_p+1)/3\ge p^4/48$, the condition
$p^2\ge128M$ proves the two stated upper bounds. Four real numbers with successive
gaps at least one have Vandermonde modulus at least $1\cdot2\cdot3\cdot1\cdot2\cdot1=12$.
At a selected root the product of its three difference moduli is at least two
and at most $S^3$. These prove the displayed discriminant and derivative bounds.

Finally $11\mid p+1$. The separately original state
\[
 R_0=11,\quad a_0=u_0=(p+11)/4,\quad
 (h,r,s,\kappa)=(a_0,1,1,(p+1)/11)
\]
returns $(a_0,a_0\kappa,pa_0\kappa)$. This proves the last assertion and excludes
any interpretation of the construction as an ES counterexample.
\end{proof}

For $\Sigma=\{2,3,5,7,11\}$ and $l=13$, the first certified example used here is
\[
 p=207481,\quad a=51871,\quad b=95637744510,
 \quad c=\frac{4960825445478210}{286913181659}.             \tag{IC35}
\]
The nonintegral denominator in the rational ES triple is
\[
 z=\frac{1029277024253264489010}{286913181659},\qquad v_{13}(z)=-1.
\]
Its independent original endpoint solution has denominators
\[
 (51873,978428526,203005329003006).
\]
The certificate also checks prime parameters $567841$ and $687961$ in the same
progression. Trial division proves each of these primality assertions.

Theorem~\ref{thm:crtcontrols} does not concern all places simultaneously, nor a
$p$-dependent complete set of possible denominator primes. The omitted prime is
an explicit place where the target fails the original integrality condition.
It shows why a fixed finite set of local checks, even with these exact real
bounds and the rank-eight completion, is not itself that condition.
